Prima di affrontare cifrari moderni, analizziamo cifrari storici per capirne le debolezze
\subsection{Cifrario a Shift}
    Questo tipo di cifrario è anche detto Cifrario di Cesare
    \begin{definition}[Cifrario a shift]
        Sia $M = C = K = \mathbb{Z}_{26}$ per $0 \leq k \leq 25$, per ogni messaggio $x \in M$
        \begin{align*}
            e_K(x) = (x + K) \mod 26 \\
            d_K(y) = (y - K) \mod 26 \\
        \end{align*}
    \end{definition}
    \begin{example}
        Sia $K = 11$ per $x = \texttt{wewillmeet}$

        \vspace{1em}
        \begin{table}[h]
            \begin{tabular}{l| *{10}{c}}
                x & 22 & 4 & 22 & 8 & 11 & 11 & 12 & 4 & 4 & 19 \\
                \hline
                y & 7 & 15 & 7 & 19 & 22 & 22 & 23 & 15 &15 & 4 
            \end{tabular}

            \vspace{1em}
            La $y$ corrispondente è \texttt{HPHTWWXPPE}.
        \end{table}
    \end{example}
    
    \noindent
    Per rompere il sistema crittografico è sufficiente provare tutte le 26 possibili chiavi, computazionalmente facile.
    
    \subsection{Cifrario a Sostituzione}
    \begin{definition}[Cifrario a sostituzione] 
        Sia $M = C = \mathbb{Z}_{26}$. $K$ è composto da tutte le possibili permutazioni dei 26 simboli $0,1, \dots, 25$\\
        $\forall \pi \in K$ definiamo
        \[ e_{\pi}(x) = \pi(x) \]
        \[ d_{\pi}(y) = \pi^{-1}(y)\]
        con $\pi^{-1}$ la permutazione inversa di $\pi$.
    \end{definition}
    \noindent
    Le possibili chiavi del cifrario sono tutte le possibili permutazioni $26!$, l'attacco per enumerazione non è ragionevole.
    
    \begin{example}
    Sia $K$ la chiave definita nel seguente modo
    \begin{table}[h]
        \begin{tabular}{c*{25}{|c}}
            a & b & c & d & e & f & g & h & i & j & k & l & m & n & o & p & q & r & s & t & u & v & w & x & y & z \\
            \hline
            X & N & Y & A & H & P & O & G & Z & Q & W & B & T & S & F & L & R & C & V & M & U & E & K & J & D & I
        \end{tabular}
    \end{table}

    \noindent
    per $y = \texttt{MGZVYZLGHCMHJMYXSSFMNHAHYCDLMMA}$, $x = \pi^{-1}(y) = \texttt{thisciphertextcannotbedecrypted}$ 
    \end{example}
    
    \subsubsection{Crittoanalisi}
    È possibile decriptare il messaggio (trovando la chiave) attraverso l'analisi delle frequenze delle lettere presenti nel \textit{ciphertext}. L'analisi delle 10 lettere più frequenti nel testo cifrato permette l'abbinamento con le lettere corrispondenti.

    Una volta stabilito un piccolo insieme di lettere di cui si ha abbastanza certezza, è possibile valutare la frequenza delle coppie di 2 lettere più comuni. Aggiungendo nuove assunzioni è possibile ottenere sempre più informazioni sulla permutazione utilizzata.
    
    \subsection{Cifrario affine}
        %Definizione 2.3
        \begin{definition}[Funzione di Eulero]
            Siano $a,m \in \mathbb{Z}. \ a\geq 1, m \geq 2$ Se $\gcd(a,m) = 1$ $a$ e $m$ si dicono \textit{co-primi} tra loro. Il numero di interi in $\mathbb{Z}_m$ co-primi con $m$ è scritto come $\phi(m)$ \textit{(funzione di Eulero)} 
        \end{definition}
    
        
        %Teorema 2.2
        \begin{theorem} Supponiamo $m$ un intero con la seguente fattorizzazione 
            \[m = \prod^{n}_{i=1}{p_i}^{e_i}\]
            dove $p$ è l'\textit{i-esimo} numero primo, e $e_i > 0$. Allora  \[\phi(m) = \prod^{n}_{i=1}{({p_i}^{e_i}-{p_i}^{e_i-1})}\]
        \end{theorem}

        \begin{definition}
            Sia $a  \in \mathbb{Z}_m$. L'inversa moltiplicativa di $a$ modulo $m$ scritta $a^{-1} \mod m$, è un elemento $a' \in \mathbb{Z}_m$ tale che $aa' \equiv 1 \pmod m$.
            L'inversa moltiplicativa esiste solo se $a$ e $m$ sono \textit{co-primi}
        \end{definition}

        \begin{definition}[Cifrario Affine]
            Sia $M = C = \mathbb{Z}_{26}$ e $K = \{ (a,b) \in \mathbb{Z}_{26} \times \mathbb{Z}_{26} \mid \gcd(a,26) = 1 \}$                
            per $K = (a,b)$ si definiscono
            
            \begin{align*}
                e_K(x) = (ax + b) \mod 26\\
                d_K(y) = a^{-1}(y-b) \mod 26
            \end{align*}
        \end{definition}

        \begin{example}
            \begin{align*}
                &K = (7,3) \quad 7^{-1} \mod 26 = 15, \quad e_K(x) = 7x + 3 \mod 26, \quad d_K(y) = 15x - 19 \mod 26\\
                &x = \texttt{hot} = 7 \ 14 \ 19 \quad e_K(x) = \underbrace{7\cdot 7 + 3}_{52} \ \underbrace{14\cdot 7 + 3}_{101} \ \underbrace{19 \cdot 7 + 3}_{136} = \texttt{AXG}
            \end{align*}
        \end{example}
    
    \subsection{Cifrario di Vigenère}
    \begin{definition}[Cifrario di Vigenère]
        Sia $m \in \mathbb{Z}$, $M = C = K = (\mathbb{Z}_{26})^{m}$. Per una chiave $K = (k_{1},\dots,k_{26})$, si definiscono per $x = (x_1,\dots,x_m) \in M$ e $y = (y_1,\dots,y_m) \in C$
        \begin{align*}
            e_K(x) = (x_1+k_1,\dots, x_m + k_m)\\    
            d_K(y) = (y_1-k_1,\dots,y_m-k_m)
        \end{align*}
    \end{definition}
    
    \begin{example}
        $m = 6$ e $K = \texttt{chipher} = (2,8,15,7,4,17)$,
        \[
            x = \texttt{thiscryptosystemisnotsecure} \qquad y = \texttt{VPXZGIAXIVWPUBTTMJ}
        \]
    \end{example}
\subsubsection{Crittoanalisi}
Per rompere il cifrario è necessario procedere per passi, il primo passo consiste nel determinare la lunghezza della chiave $m$, attraverso il test di \emph{Kasiski} e l'indice di coincidenza.

\begin{enumerate}
    \item Il test di \emph{Kasiski} consiste nell'osservare due segmenti identici, nel \emph{plaintext} saranno criptati nella stessa forma se sono a $\delta$ posizioni di distanza con $\delta \equiv 0 \mod m$
    
    \item Indice di coincidenza, per una stringa $x=(x_1,\dots,x_m)$ l'indice di coincidenza $I_c(x)$ misura la probabilità che due elementi casuali in $x$ siano uguali
    
    \begin{align*}
        I_c(x) = \frac{\sum_{i=0}^{25}\binom{f_i}{2}}{\binom{n}{2}} = \frac{\sum_{i=0}^{25} f_i(f_i-1)}{n(n-1)}    
    \end{align*}
    dove $f_i$ è la frequenza del carattere \emph{i-esimo}.
    Per un testo in lingua inglese ci si aspetta che

    $I_c(x) = \sum_{i=0}^{25}p_{i}^{2} = 0.0065$. È sufficiente verificare che 
    \[
        M_g = \sum_{i=0}^{25} \frac{p_i f_{i+g}}{n/m} 
    \]
    Si definisce una sequenza di stringhe $y$ tale che
    \begin{align*}
        y_1 &= y_1\, y_{m+1}\, y_{2m + 1}\\
        y_2 &= y_2\, y_{m+2}\, y_{2m + 2}\\
        y_m &= y_m\, y_{2m} \, y_{3m}\\
    \end{align*}
\end{enumerate}
e si cerca il valore di $m$ per cui le sotto-stringhe $(y1,\dots,y_m)$ hanno $I_c(y_i)$ più vicino a 0.065, per ogni i = $1,\dots,m$.

Successivamente si provano 25 shift per le sotto-stringhe con $M_g = \frac{p_i f_{i+g}}{n/m}$ (si cerca di accoppiare le frequenze delle lettere nel testo con $p_i$, frequenza delle lettere nella lingua inglese), il valore più vicino a 0.065 ci permette di determinare il carattere in $K$.
% 1-2-3-4
% Cifrario di Viginère, Cifrario di Cesare, Cifrario per Sostituzione, Cifrario di Hill

\subsection{Cifrario di Hill}
\begin{definition}[Cifrario di Hill]
    Sia $m \geq 2 \in \mathbb{Z}$, $M = C = (\mathbb{Z})^{26}$ e $K = \{ \text{matrici invertibili } m \times m\}$, si definiscono 
    \[
        e_K(x) = xK \qquad d_K(y) = yK^{-1}
    \]
\end{definition}
\begin{example}
Sia $K = \begin{pmatrix} 11 & 8 \\ 3 & 7\end{pmatrix}$, $K^{-1} = \begin{pmatrix} 7 & 18\\ 23 & 11 \end{pmatrix}$, 
\[
    x = \texttt{july} = \begin{pmatrix}9 & 20\end{pmatrix},\begin{pmatrix}11 & 24\end{pmatrix} 
\]
\begin{align*}
    \begin{pmatrix}9 & 20\end{pmatrix}K &= (3,4)\\
    \begin{pmatrix}11 & 24\end{pmatrix}K &= (11,22)
\end{align*}
$y = \texttt{DELW}$.
\end{example}

\subsubsection{Crittoanalisi}
Il cifrario è facile da rompere con un \emph{known-plaintext attack}. Supponendo che l'avversario abbia $m$ coppie $x,y$,$x \in M, y \in C$ tale che $e_K(x) = y$
